% Part: normal-modal-logic
% Chapter: filtrations
% Section: introduction

\documentclass[../../../include/open-logic-section]{subfiles}

\begin{document}

\olfileid{nml}{fil}{int}

\olsection{مقدمه}

یکی از پرسش‌های مهم دربارهٔ هر منطق همواره این است که آیا تصمیم‌پذیر
است؛ یعنی آیا روشی مؤثر وجود دارد که به پرسش «آیا این !!{formula}
معتبر است؟» پاسخ دهد. منطق گزاره‌ای تصمیم‌پذیر است: می‌توانیم با ساختن
جدول ارزش به‌طور مؤثر بیازماییم که آیا !!a{formula} همان‌گویی است، و
برای یک !!{formula} مفروض، جدول ارزش متناهی است. اما آشکارا نمی‌توانیم
بیازماییم که آیا یک فرمول موجهاتی در همهٔ مدل‌ها صادق است، زیرا
بی‌نهایت مدل وجود دارد. می‌توانیم همهٔ مدل‌های متناهیِ مربوط به یک
فرمول مفروض را فهرست کنیم، زیرا تنها تخصیص زیرمجموعه‌های جهان‌ها به
!!{propositional variable}هایی اهمیت دارد که واقعاً در !!{formula}
ظاهر می‌شوند. اگر رابطهٔ دسترسی‌پذیری ثابت باشد، تخصیص‌های گوناگونِ
ممکن~$V(p)$ صرفاً همهٔ زیرمجموعه‌های~$W$ هستند، و اگر~$\card{W} = n$،
تعدادشان~$2^n$ است. اگر !!{formula}~$!A$ شامل~$m$
!!{propositional variable} باشد، در این صورت~$2^{nm}$ مدل گوناگون با
$n$ جهان وجود دارد. برای هریک می‌توانیم صرفاً با محاسبهٔ ارزش صدق~$!A$
در هر جهان بیازماییم که آیا~$!A$ در همهٔ جهان‌ها صادق است. البته باید
همهٔ روابط دسترسی‌پذیری ممکن را نیز بررسی کنیم، اما شمار روابط روی
$n$ جهان نیز متناهی است (به‌طور مشخص، شمار زیرمجموعه‌های~$W \times W$،
یعنی~$2^{n^2}$).

اگر به منطق~\Log{K} علاقه‌مند نباشیم و منطق مورد نظر به‌وسیلهٔ رده‌ای
از مدل‌ها تعریف شده باشد (برای نمونه، مدل‌های بازتابی و متعدی)، باید
بتوانیم بیازماییم که آیا رابطهٔ دسترسی‌پذیری از نوع درست است. هرگاه
قاب‌های مورد نظر ما با !!{formula}های موجهاتی تعریف‌پذیر باشند، این
کار ممکن است (برای نمونه، با آزمودن اینکه آیا~\Ax{T} و~\Ax{4} در قاب
معتبرند). پس ایده این است که همهٔ قاب‌های متناهی را یکی‌یکی بررسی
کنیم، برای هرکدام بیازماییم که آیا عضوی از ردهٔ مورد نظر است، سپس همهٔ
مدل‌های ممکن روی آن قاب را فهرست کنیم و بیازماییم که آیا~$!A$ در هریک
صادق است. اگر نبود، متوقف می‌شویم: $!A$ در ردهٔ مدل‌های مورد نظر معتبر
نیست.

این ایده مشکلی دارد: نمی‌دانیم چه هنگام، اگر اصلاً زمانی باشد، می‌توانیم
جست‌وجو را متوقف کنیم. اگر !!{formula} یک پادمدل متناهی داشته باشد،
روش ما آن را خواهد یافت. اما اگر هیچ پادمدل متناهی نداشته باشد، پاسخی
به دست نمی‌آوریم. ممکن است !!{formula} معتبر باشد (یعنی اصلاً پادمدلی
نداشته باشد)، یا تنها پادمدلی نامتناهی داشته باشد که هرگز آن را بررسی
نخواهیم کرد. اگر بتوانیم نشان دهیم هر !!{formula} دارای پادمدل، یک
پادمدل متناهی نیز دارد، می‌توان بر این مشکل غلبه کرد. در این صورت
می‌گوییم منطق دارای \emph{خاصیت مدل متناهی} است.

اما چگونه نشان می‌دهیم یک منطق خاصیت مدل متناهی دارد؟ یک راه این است
که روشی برای تبدیل یک (پاد)مدل نامتناهیِ~$!A$ به مدلی متناهی بیابیم.
اگر این کار شدنی باشد، هرگاه مدلی وجود داشته باشد که~$!A$ در آن صادق
نیست، مدل متناهیِ حاصل نیز~$!A$ را کاذب می‌کند. آن مدل متناهی در فهرست
همهٔ مدل‌های متناهی ما ظاهر خواهد شد و سرانجام برای هر فرمول نامعتبر
تعیین خواهیم کرد که نامعتبر است. اگر فرمول معتبر باشد، روش ما پایان
نمی‌یابد. اگر افزون‌براین بتوانیم نشان دهیم اندازهٔ مدل متناهی‌ای که
روش ما فراهم می‌کند کران بالایی دارد و این کران تنها به
!!{formula}~$!A$ وابسته است، آنگاه اندازه‌ای خواهیم داشت که در جست‌وجوی
پادمدل باید مدل‌های متناهی را تنها تا آن اندازه بیازماییم. اگر تا آن
زمان پادمدلی نیافته باشیم، هیچ پادمدلی وجود ندارد. در این صورت روش ما
واقعاً برای هر !!{formula}~$!A$ دربارهٔ پرسش «آیا~$!A$ معتبر است؟»
تصمیم خواهد گرفت.

راهبردی که اغلب برای تبدیل ساختارهای نامتناهی به ساختارهای متناهی سودمند
است، «یکی‌گرفتن» !!{element}هایی از ساختار است که از جهات مربوط رفتاری
یکسان دارند. اگر بی‌نهایت جهان در~$\mModel{M}$ از جهات مربوط رفتاری
یکسان داشته باشند، ممکن است بتوانیم \emph{همهٔ} جهان‌ها را در شمار
متناهی‌ای از «رده‌های» چنین جهان‌هایی گرد آوریم؛ هریک از این رده‌ها
می‌تواند متناهی یا نامتناهی باشد. به بیان دیگر، باید مجموعهٔ جهان‌ها را
به شیوهٔ مناسبی افراز کنیم: افرازی با تنها شمار متناهی بلوک که هریک
می‌تواند متناهی یا نامتناهی باشد. سپس مدل جدید~$\mModel{M^*}$ را چنان
تعریف می‌کنیم که جهان‌هایش همین بلوک‌ها باشند. شمار متناهی بلوک در مدل
قدیم، شمار متناهی جهان در مدل جدید، یعنی مدلی متناهی، به دست می‌دهد.
بلوک شامل جهان~$w$ را~$[w]$ می‌نامیم. می‌خواهیم
$\mSat{M}{!A}[w]$ اگر و تنها اگر~$\mSat{M^*}{!A}[{[w]}]$؛ زیرا
می‌خواهیم اگر مدل قدیم پادمدلی برای~$!A$ بود، مدل جدید نیز چنین باشد.
این امر مستلزم آن است که هم افراز و هم رابطهٔ دسترسی‌پذیریِ
$\mModel{M^*}$ را به شیوهٔ مناسب تعریف کنیم.

برای دیدن چگونگی کار، نخست فرض کنید رابطهٔ دسترسی‌پذیری همگانی است.
$\mSat{M}{\Box !B}[w]$ اگر و تنها اگر برای هر~$v \in W$،
$\mSat{M}{!B}[v]$؛ و برای~$\mModel{M^*}$ نیز همین‌گونه است، جز آنکه
$[w]$ و~$[v]$ را داریم. به‌عنوان نخستین ایده، می‌گوییم دو جهان~$u$ و
$v$ هم‌ارزند (به یک بلوک تعلق دارند) اگر دربارهٔ همهٔ
!!{propositional variable}ها در~$\mModel{M}$ توافق داشته باشند؛ یعنی
$\mSat{M}{p}[u]$ اگر و تنها اگر~$\mSat{M}{p}[v]$. بگذارید
$V^*(p) = \Setabs{[w]}{\mSat{M}{p}[w]}$. هدف ما اثبات این است که
$\mSat{M}{!A}[w]$ اگر و تنها اگر~$\mSat{M^*}{!A}[{[w]}]$. آشکارا
این امر را با استقرا اثبات می‌کنیم: حالت پایه~$!A \equiv p$ است. نخست
فرض کنید~$\mSat{M}{p}[w]$. آنگاه بنا به تعریف~$[w] \in V^*$، پس
$\mSat{M^*}{p}[{[w]}]$. اکنون فرض کنید
$\mSat{M^*}{p}[{[w]}]$. این یعنی~$[w] \in V^*(p)$؛ یعنی برای یک
$v$ هم‌ارز با~$w$، $\mSat{M}{p}[v]$. اما «$w$ با~$v$ هم‌ارز است»
یعنی «$w$ و~$v$ دقیقاً همان !!{propositional variable}ها را صادق
می‌کنند»؛ پس~$\mSat{M}{p}[w]$. اکنون برای گام استقرا، برای نمونه
$!A \ident \lnot !B$. آنگاه~$\mSat{M}{\lnot !B}[w]$ اگر و تنها اگر
$\mSat/{M}{!B}[w]$، اگر و تنها اگر
$\mSat/{M^*}{!B}[{[w]}]$ (بنا بر فرض استقرا)، اگر و تنها اگر
$\mSat{M^*}{\lnot !B}[{[w]}]$. برای عملگرهای غیرموجهاتی دیگر نیز
به‌همین‌سان است. برای~$\Box$ نیز کار می‌کند: فرض کنید
$\mSat{M^*}{\Box !B}[{[w]}]$. این یعنی برای هر~$[u]$،
$\mSat{M^*}{!B}[{[u]}]$. بنا بر فرض استقرا، برای هر~$u$،
$\mSat{M}{!B}[u]$. در نتیجه~$\mSat{M}{\Box !B}[w]$.

در حالت کلی که باید رابطهٔ دسترسی‌پذیریِ~$\mModel{M^*}$ را نیز تعریف
کنیم، امور پیچیده‌ترند. مدل~$\mModel{M^*}$ را \emph{تصفیه} می‌نامیم
اگر رابطهٔ دسترسی‌پذیری~$R^*$ آن شرایط لازم برای پیش‌رفتن اثبات
استقرایی بالا را برآورده کند. در این صورت هر تصفیه~$\mModel{M^*}$،
$!A$ را در~$[w]$ صادق می‌کند اگر و تنها اگر~$\mModel{M}$، $!A$ را
در~$w$ صادق کند. بااین‌حال، اکنون باید نشان دهیم تصفیه‌ها \emph{وجود
دارند}؛ یعنی می‌توانیم~$R^*$ را چنان تعریف کنیم که شرایط لازم را
برآورده سازد. برای شدنی‌بودن این کار باید جهان‌های~$u$ و~$v$ را نه
فقط هنگامی هم‌ارز بشماریم که دربارهٔ همهٔ !!{propositional variable}ها
توافق دارند، بلکه باید دربارهٔ همهٔ زیر-!!{formula}های~$!A$ نیز توافق
داشته باشند. چون~$!A$ تنها شمار متناهی‌ای زیر-!!{formula} دارد، این
امر همچنان تضمین می‌کند که تصفیه متناهی باشد. تعریف تصفیه تنها یک راه
ندارد؛ و برای اطمینان از اینکه رابطهٔ دسترسی‌پذیریِ تصفیه خواص لازم
(برای نمونه، بازتابی، متعدی و مانند آن) را دارد، باید در تعریف~$R^*$
ابتکار به خرج دهیم.


\end{document}
